% SGA 3, Expose I, Loop 1 English reconstruction.
% Scope: Section 1.8 through the end of Section 2.2, including the transition
% sentence immediately before Section 2.3.
% Authority: Polo--Gille Exp1-13oct24.pdf, component pages 10--15
% (combined-reader pages 20--25). OCR and the Reinhold English Markdown are
% drafting/comparison witnesses only.

\SGARefTarget{sga3:i:subsection:1:8}
\subsection*{1.8. 常对象}

设 \(\mathscr C\) 是一个范畴, 其中存在直和与纤维积, 且直和与基变换可交换
(例如概形范畴).\footnote{编者注: 旧术语“预概形/概形”已替换为现行术语
“概形/分离概形”.} 对每个集合 \(E\) 和 \(\mathscr C\) 的每个对象 \(S\), 置
\begin{equation}
  \tag{1.8.1}\label{sga3:i:subsection:1-8:equation:1-8-1}
  E_S=
  \left\{
    \begin{aligned}
      &\text{对象族 }(S_i)_{i\in E}\text{ 的直和,}\\
      &\text{其中各对象属于 }\mathscr C\text{ 且均同构于 }S.
    \end{aligned}
  \right.
\end{equation}
这个对象由下式刻画:
\begin{equation}
  \tag{1.8.2}\label{sga3:i:subsection:1-8:equation:1-8-2}
  \operatorname{Hom}_{\mathscr C}(E_S,T)
  =\operatorname{Hom}_{\mathbf{Set}}
      \bigl(E,\operatorname{Hom}_{\mathscr C}(S,T)\bigr),
\end{equation}
其中 \(T\in\operatorname{Ob}(\mathscr C)\) 任意.

对象 \(E_S\) 有一个到 \(S\) 的典范投影, 因而 \(E\mapsto E_S\) 实际上是从
\(\mathbf{Set}\) 到 \(\mathscr C_{/S}\) 的函子.\footnote{编者注: 以下三段调整了
语句次序, 并增补了关于下述条件 \((*)\) 所起作用的若干细节.} 若 \(S'\to S\) 是
\(\mathscr C\) 中的箭头, 则由于直和与基变换可交换, 有
\[
  E_{S'}=(E_S)_{S'}.
\]
特别地, 若 \(\mathscr C\) 有终对象 \(e\), 则
\[
  E_S=(E_e)_S.
\]

从 \(\mathbf{Set}\) 到 \(\mathscr C_{/S}\) 的函子 \(E\mapsto E_S/S\) 与有限积
可交换. 只需证明
\begin{equation}
  \tag{\ensuremath{\times}}
  E_S\mathbin{\times_S}F_S=(E\times F)_S.
\end{equation}
事实上, 将 \SGARefLink{sga3:i:subsection:1:7}{1.7} 的结果用于 \(\mathscr C_{/S}\),
可知对每个 \(T\in\operatorname{Ob}(\mathscr C_{/S})\) 都有自然同构
(所有未注明下标的 \(\operatorname{Hom}\) 均在 \(\mathscr C_{/S}\) 中取):
\begin{align*}
  \operatorname{Hom}((E\times F)_S,T)
  &\simeq
    \operatorname{Hom}_{\mathbf{Set}}
      \bigl(E\times F,\operatorname{Hom}(S,T)\bigr)\\
  &\simeq
    \operatorname{Hom}_{\mathbf{Set}}
      \bigl(E,
        \operatorname{Hom}_{\mathbf{Set}}
          (F,\operatorname{Hom}(S,T))\bigr)\\
  &\simeq
    \operatorname{Hom}_{\mathbf{Set}}
      \bigl(E,\operatorname{Hom}(F_S,T)\bigr),
\end{align*}
以及
\begin{align*}
  \operatorname{Hom}(E_S\mathbin{\times_S}F_S,T)
  &\simeq
    \operatorname{Hom}
      \bigl(E_S,\underline{\operatorname{Hom}}(F_S,T)\bigr)\\
  &\simeq
    \operatorname{Hom}_{\mathbf{Set}}
      \bigl(E,
        \operatorname{Hom}
          (S,\underline{\operatorname{Hom}}(F_S,T))\bigr).
\end{align*}
但
\[
  \operatorname{Hom}
    \bigl(S,\underline{\operatorname{Hom}}(F_S,T)\bigr)
  \simeq \operatorname{Hom}(F_S,T),
\]
这就证明了 \((\times)\).

设以 \(\varnothing\) 表示 \(\mathscr C\) 的一个始对象, 且下图是笛卡尔图:

% Authority: Exp1-13oct24.pdf p. 11; combined reader p. 21.
\begin{center}
  \makebox[0.62\textwidth][l]{%
    \textup{(*)}\hfill
    \begin{tikzcd}[column sep=4.2em,row sep=3.0em]
  \varnothing
    \arrow[r]
    \arrow[d]
  &
  S
    \arrow[d]
  \\
  S
    \arrow[r]
  &
  S\amalg S
\end{tikzcd}
%
    \hfill\phantom{\textup{(*)}}%
  }
\end{center}
概形范畴满足这一条件.\footnote{编者注: 若 \(\mathscr C\) 是仅有箭头
\(A\to B\), \(\operatorname{id}_A\) 和 \(\operatorname{id}_B\) 的范畴, 则条件
\((*)\) 不成立. 此时有 \(B\coprod B=B\) 和 \(B\times_BB=B\not\simeq A\).}
那么函子 \(E\mapsto E_S\) 与有限逆极限可交换.

事实上, 由 \((\times)\) 可知, 只需证明 \(E\mapsto E_S\) 与纤维积可交换.
设 \(u:E\to G\) 和 \(v:F\to G\) 是两个集合映射. 由于 \(\mathscr C\) 中的直和
与基变换可交换, 有
\begin{equation}
  \tag{1}\label{sga3:i:subsection:1-8:equation:1}
  F_S\mathbin{\times_{G_S}}E_S
  \simeq
  \coprod_{f\in F}
    \bigl(S_f\mathbin{\times_{G_S}}E_S\bigr)
  \simeq
  \coprod_{\substack{f\in F\\x\in E}}
    \bigl(S_f\mathbin{\times_{G_S}}S_x\bigr).
\end{equation}
若 \(v(f)\ne u(x)\), 则在 \(\mathscr C_{/S}\) 中有态射
\[
  S_f\mathbin{\times_{G_S}}S_x
  \longrightarrow
  S_f\mathbin{\times_{S_{v(f)}\coprod S_{u(x)}}}S_x.
\]
由条件 \((*)\), 右端为 \(\varnothing\). 因此
\begin{equation}
  \tag{2}\label{sga3:i:subsection:1-8:equation:2}
  S_f\mathbin{\times_{G_S}}S_x\simeq\varnothing
  \qquad\text{当 }v(f)\ne u(x).
\end{equation}
另一方面, 若 \(v(f)=u(x)\), 则在 \(\mathscr C_{/S}\) 中有态射
\[
  S\longrightarrow S_f\mathbin{\times_{G_S}}S_x.
\]
由于 \(S\xrightarrow{\operatorname{id}}S\) 是 \(\mathscr C_{/S}\) 的终对象, 可得
\begin{equation}
  \tag{3}\label{sga3:i:subsection:1-8:equation:3}
  S_f\mathbin{\times_{G_S}}S_x\simeq S
  \qquad\text{当 }v(f)=u(x).
\end{equation}
结合 \SGARefLink{sga3:i:subsection:1-8:equation:1}{(1)}, \SGARefLink{sga3:i:subsection:1-8:equation:2}{(2)}
及 \SGARefLink{sga3:i:subsection:1-8:equation:3}{(3)}, 得到关于 \(S\) 函子性的同构:
\[
  F_S\mathbin{\times_{G_S}}E_S
  \simeq \coprod_{F\times_G E}S
  =(F\mathbin{\times_G}E)_S.
\]

形如 \(E_S\) 的对象称为\emph{常对象 (constant object)}. 存在一个关于 \(E\)
函子性的态射
\[
  E\longrightarrow\Gamma(E_S/S),
\]
它把每个 \(i\in E\) 对应到由 \(S\) 到 \(S_i\) 的同构所定义的 \(E_S\) 在
\(S\) 上的截面. 设条件 \((*)\) 对 \(\mathscr C\) 的每个对象 \(S\) 均成立.
那么只要 \(S\not\simeq\varnothing\), \(E\to\Gamma(E_S/S)\) 就是单态射.

若 \(\mathscr C\) 是概形范畴, 则 \(\Gamma(E_S/S)\) 等同于从拓扑空间 \(S\)
到集合 \(E\) 的局部常值映射的集合; 上述映射把 \(E\) 的每个元素对应到相应的
常值映射. 因此, 也可以等价地将 \(E_S\) 定义为下述函子的代表对象: 该函子把每个
\(S\) 上的 \(S'\) 对应到从拓扑空间 \(S'\) 到 \(E\) 的局部常值函数的集合.
\footnote{编者注: 对角态射 \(E_S\to E_S\times_SE_S=(E\times E)_S\) 也是闭浸入;
换言之, \(E_S\) 在 \(S\) 上分离.}

\section*{\SGAInlineRefTarget{sga3:i:section:2}2. 代数结构}
\addcontentsline{toc}{section}{2. 代数结构}

给定集合范畴中的一种代数结构, 我们要将它推广到范畴 \(\mathscr C\).
先讨论群的例子.

\subsection*{\SGAInlineRefTarget{sga3:i:subsection:2-1}2.1. 群结构}

沿用上一段的记号.

\medskip
\noindent\SGAInlineRefTarget{sga3:i:definition:2-1-1}\textbf{定义 2.1.1.}
设 \(\mathbf G\in\operatorname{Ob}(\widehat{\mathscr C})\).
\(\mathbf G\) 上的一个\emph{\(\widehat{\mathscr C}\)-群结构}是指: 对每个
\(S\in\operatorname{Ob}(\mathscr C)\), 给定集合 \(\mathbf G(S)\) 上的一个群结构,
使得对 \(\mathscr C\) 中的每个箭头 \(f:S'\to S''\), 映射
\[
  \mathbf G(f):\mathbf G(S'')\longrightarrow\mathbf G(S')
\]
是群同态. 若 \(\mathbf G\) 与 \(\mathbf H\) 是两个 \(\widehat{\mathscr C}\)-群,
则从 \(\mathbf G\) 到 \(\mathbf H\) 的\emph{\(\widehat{\mathscr C}\)-群态射}是指
满足下述条件的态射 \(u\in\operatorname{Hom}(\mathbf G,\mathbf H)\): 对每个
\(S\in\operatorname{Ob}(\mathscr C)\), \(u(S):\mathbf G(S)\to\mathbf H(S)\)
都是群同态.

以
\(\operatorname{Hom}_{\widehat{\mathscr C}\text{-}\mathrm{Gr}}
  (\mathbf G,\mathbf H)\)
表示这些态射的集合, 并以 \((\widehat{\mathscr C}\text{-}\mathrm{Gr})\) 表示它们构成的范畴.

\medskip
\noindent\textbf{例.}
对 \(\mathbf E\in\operatorname{Ob}(\widehat{\mathscr C})\), 对象
\(\underline{\operatorname{Aut}}(\mathbf E)\) 显然具有
\(\widehat{\mathscr C}\)-群结构. 终对象 \(\underline{\mathbf e}\) 有唯一的
\(\widehat{\mathscr C}\)-群结构, 使它成为 \((\widehat{\mathscr C}\text{-}\mathrm{Gr})\)
的终对象.

对每个 \(S\in\operatorname{Ob}(\mathscr C)\), 设 \(e_{\mathbf G}(S)\) 为
\(\mathbf G(S)\) 的单位元. 这些元素 \(e_{\mathbf G}(S)\) 组成的族定义了
\[
  e_{\mathbf G}\in\Gamma(\mathbf G)
  =\operatorname{Hom}(\underline{\mathbf e},\mathbf G).
\]
它是一个 \(\widehat{\mathscr C}\)-群态射 \(\underline{\mathbf e}\to\mathbf G\),
称为 \(\mathbf G\) 的\emph{单位截面 (unit section)}.

给定 \(\mathbf G\) 上的 \(\widehat{\mathscr C}\)-群结构, 等价于给定
\(\mathbf G\) 上的一个合成律, 即一个 \(\widehat{\mathscr C}\)-态射
\[
  \pi_{\mathbf G}:\mathbf G\times\mathbf G\longrightarrow\mathbf G
\]
使得对每个 \(S\in\operatorname{Ob}(\mathscr C)\), \(\pi_{\mathbf G}(S)\) 都赋予
\(\mathbf G(S)\) 一个群结构.

同样, \(f:\mathbf G\to\mathbf H\) 是 \(\widehat{\mathscr C}\)-群态射, 当且仅当下图交换:

% Authority: Exp1-13oct24.pdf p. 13; combined reader p. 23.
\begin{center}
  \begin{tikzcd}[column sep=4.5em,row sep=3.0em]
  \mathbf G\times\mathbf G
    \arrow[r,"\pi_{\mathbf G}"]
    \arrow[d,"{(f,f)}"']
  &
  \mathbf G
    \arrow[d,"f"]
  \\
  \mathbf H\times\mathbf H
    \arrow[r,"\pi_{\mathbf H}"]
  &
  \mathbf H
  \qquad .
\end{tikzcd}

\end{center}

若 \(\mathbf G\) 的子对象 \(\mathbf H\) 满足: 对每个 \(S\), \(\mathbf H(S)\)
均为 \(\mathbf G(S)\) 的子群, 则它显然具有由 \(\mathbf G\) 的结构诱导出的
\(\widehat{\mathscr C}\)-群结构. 这是使单态射 \(\mathbf H\to\mathbf G\) 成为
\(\widehat{\mathscr C}\)-群态射的唯一结构. 配备这一结构后, 称 \(\mathbf H\)
为 \(\mathbf G\) 的一个\emph{子 \(\widehat{\mathscr C}\)-群}.

若 \(\mathbf G\) 与 \(\mathbf H\) 是两个 \(\widehat{\mathscr C}\)-群, 则
\(\mathbf G\times\mathbf H\) 具有显然的积结构: 对每个 \(S\), 在集合
\(\mathbf G(S)\times\mathbf H(S)\) 上赋予所给群结构的积. 所得对象称为
\(\mathbf G\) 与 \(\mathbf H\) 的\emph{积 \(\widehat{\mathscr C}\)-群};
它是它们在 \(\widehat{\mathscr C}\)-群范畴中的积.

若 \(\mathbf G\) 是 \(\widehat{\mathscr C}\)-群, 则对每个
\(S\in\operatorname{Ob}(\mathscr C)\), \(\mathbf G_S\) 都是
\(\widehat{\mathscr C_{/S}}\)-群. 对两个 \(\widehat{\mathscr C}\)-群
\(\mathbf G\) 和 \(\mathbf H\), 以下式定义 \(\widehat{\mathscr C}\) 的对象
\(\underline{\operatorname{Hom}}_{
  \widehat{\mathscr C}\text{-}\mathrm{Gr}}(\mathbf G,\mathbf H)\):
\[
  \underline{\operatorname{Hom}}_{
    \widehat{\mathscr C}\text{-}\mathrm{Gr}}(\mathbf G,\mathbf H)(S)
  =
  \operatorname{Hom}_{
    \widehat{\mathscr C_{/S}}\text{-}\mathrm{Gr}}
    (\mathbf G_S,\mathbf H_S).
\]
注意, \(\underline{\operatorname{Hom}}_{
  \widehat{\mathscr C}\text{-}\mathrm{Gr}}\) 一般不是
\(\widehat{\mathscr C}\)-群, 更不是范畴 \((\widehat{\mathscr C}\text{-}\mathrm{Gr})\)
中的内部 \(\underline{\operatorname{Hom}}\) 对象.

类似地定义
\[
  \underline{\operatorname{Isom}}_{
    \widehat{\mathscr C}\text{-}\mathrm{Gr}}(\mathbf G,\mathbf H),
  \qquad
  \underline{\operatorname{End}}_{
    \widehat{\mathscr C}\text{-}\mathrm{Gr}}(\mathbf G),
  \qquad
  \underline{\operatorname{Aut}}_{
    \widehat{\mathscr C}\text{-}\mathrm{Gr}}(\mathbf G).
\]

\medskip
\noindent\SGAInlineRefTarget{sga3:i:definition:2-1-2}\textbf{定义 2.1.2.}
设 \(G\in\operatorname{Ob}(\mathscr C)\). \(G\) 上的一个
\emph{\(\mathscr C\)-群结构}是指
\(\mathbf h_G\in\operatorname{Ob}(\widehat{\mathscr C})\) 上的
\(\widehat{\mathscr C}\)-群结构. 从 \(\mathscr C\)-群 \(G\) 到
\(\mathscr C\)-群 \(H\) 的态射是指一个元素
\[
  u\in\operatorname{Hom}(G,H)
  \simeq\operatorname{Hom}(\mathbf h_G,\mathbf h_H)
\]
它定义了从 \(\mathbf h_G\) 到 \(\mathbf h_H\) 的
\(\widehat{\mathscr C}\)-群态射.

以 \((\mathscr C\text{-}\mathrm{Gr})\) 表示 \(\mathscr C\)-群的范畴.
在 \(\mathbf{Cat}\) 中有一个笛卡尔方块:

% Authority: Exp1-13oct24.pdf p. 13; combined reader p. 23.
\begin{center}
  \begin{tikzcd}[column sep=5.0em,row sep=3.2em]
  (\mathscr C\text{-}\mathrm{Gr.})
    \arrow[r]
    \arrow[d]
  &
  (\widehat{\mathscr C}\text{-}\mathrm{Gr.})
    \arrow[d]
  \\
  \mathscr C
    \arrow[r,"\mathbf h"]
  &
  \widehat{\mathscr C}
  \qquad .
\end{tikzcd}

\end{center}

因此, 只要上述定义和构造所用的函子 (积, 内部
\(\underline{\operatorname{Hom}}\) 对象等等) 可表, 这些定义和构造就立即适用于
\((\mathscr C\text{-}\mathrm{Gr})\). 它们也适用于范畴 \(\mathscr C_{/S}\).
此时将 \(\underline{\operatorname{Hom}}_{\mathscr C_{/S}\text{-}\mathrm{Gr}}\)
记为 \(\underline{\operatorname{Hom}}_{S\text{-}\mathrm{Gr}}\), 其他对象也采用类似记号.

\subsection*{\SGAInlineRefTarget{sga3:i:subsection:2-2}2.2.}

更一般地, 设 \((T)\) 是由有限逆极限定义的、以 \(n\) 个集合为基集的一种结构
---例如用笛卡尔积构成的交换图来定义的结构, 如幺半群、群、带算子的集合、环上的模
或环上的 Lie 代数结构. 上述构造定义了 \(\widehat{\mathscr C}\) 的 \(n\) 个对象
\(\mathbf F_1,\ldots,\mathbf F_n\) 上的“\((T)\) 类结构”的概念.
这种结构是指: 对每个 \(S\in\operatorname{Ob}(\mathscr C)\), 在集合
\(\mathbf F_1(S),\ldots,\mathbf F_n(S)\) 上给定一个 \((T)\) 类结构,
使得对 \(\mathscr C\) 中的每个箭头 \(S'\to S''\), 映射族
\[
  \bigl(\mathbf F_i(S'')\bigr)_i
  \longrightarrow
  \bigl(\mathbf F_i(S')\bigr)_i
\]
是关于结构种类 \((T)\) 的多重同态 (polyhomomorphism).

类似地定义 \((T)\) 类结构的态射, 得到一个范畴
\[
  (\widehat{\mathscr C}\times\cdots\times\widehat{\mathscr C})^{(T)}.
\]
于是全忠实函子 \(\mathbf h\times\cdots\times\mathbf h\) 通过逆像定义了范畴
\[
  (\mathscr C\times\cdots\times\mathscr C)^{(T)}.
\]
由于它与逆极限可交换, 在 \(\widehat{\mathscr C}\) 中引入的所有函子性性质、概念和
记号都可以移用于这一范畴.

现设 \(\mathscr C\) 中存在纤维积, 并设 \((T)\) 是这样一种代数结构:
它由笛卡尔积之间的某些态射指定, 所满足的公理由这些态射构成的交换图给出.
于是, \(\mathscr C\) 中一族对象上的一个 \((T)\) 类结构, 就由笛卡尔积之间满足
相应交换性条件的相应态射来定义. 因此, 若 \(\mathscr C\) 与 \(\mathscr C'\)
是有积的范畴, \(f:\mathscr C\to\mathscr C'\) 是与积可交换的函子, 则
\(\mathscr C\) 的一族对象 \((F_i)\) 上的 \((T)\) 类结构, 在
\(\mathscr C'\) 的对象族 \((f(F_i))\) 上诱导出一个同类结构. 这样, 每个
\(\mathscr C\)-群都成为一个 \(\mathscr C'\)-群, 由一个 \(\mathscr C\)-环和该环上的
一个模组成的每一对对象, 都成为 \(\mathscr C'\) 中相应的一对对象, 依此类推.

特别地, 设 \(\mathscr C\) 满足 \SGARefLink{sga3:i:subsection:1:8}{1.8} 的条件.
\footnote{编者注: 包括条件 \((*)\).} 该处定义的函子 \(E\mapsto E_S\)
与有限逆极限可交换. 因此, 它把群变成 \(S\)-群 (即 \(\mathscr C_{/S}\)-群),
把环变成 \(S\)-环, 等等.

\medskip
\noindent\textbf{注.}
应当注意, 将上述构造用于 \(\widehat{\mathscr C}\), 得到的恰好是该范畴中已经定义的
那些概念. 换言之, 在 \(\widehat{\mathscr C}\) 的一个对象上, 将它视为
\(\mathscr C\) 上的函子而给定 \((T)\) 类结构, 等价于在这个对象所定义的
\(\widehat{\mathscr C}\) 上的可表函子上给定 \((T)\) 类结构.
\footnote{编者注: 例如, 对于群结构, 设
\(\mathbf G\in\operatorname{Ob}(\widehat{\mathscr C})\). 若函子
\[
  \widehat{\mathscr C}\longrightarrow\mathbf{Set},
  \qquad
  \mathbf F\longmapsto
  \operatorname{Hom}_{\widehat{\mathscr C}}(\mathbf F,\mathbf G)
\]
具有群结构, 则它在 \(\mathscr C\) 上的限制, 即 \(X\mapsto\mathbf G(X)\),
也具有群结构. 反之, 若 \(\mathbf G\) 是 \(\widehat{\mathscr C}\)-群, 则其乘法态射
\(\pi_{\mathbf G}:\mathbf G\times\mathbf G\to\mathbf G\) 对每个
\(\mathbf F\in\operatorname{Ob}(\widehat{\mathscr C})\), 在
\(\operatorname{Hom}_{\widehat{\mathscr C}}(\mathbf F,\mathbf G)\) 上诱导出一个群结构,
而且这关于 \(\mathbf F\) 是函子性的.}

我们还要讨论这一构造的另外两个特殊情形: 带算子群的结构与模.

% End of component. Section 2.3 begins in the next unit.
